\documentclass[12pt,a4paper]{article}

% ── Packages ──────────────────────────────────────────────────────────
\usepackage[margin=1in]{geometry}
\usepackage{graphicx}
\usepackage{booktabs}
\usepackage{enumitem}
\usepackage{hyperref}
\usepackage{titlesec}
\usepackage{fancyhdr}
\usepackage{xcolor}
\usepackage{tabularx}
\usepackage{parskip}
\\usepackage[expansion=false]{microtype}
\usepackage{amsmath}
\usepackage{textcomp}
\usepackage[utf8]{inputenc}
\usepackage[T1]{fontenc}
\newcommand{\rupee}{\textrm{\textsf{\textup{Rs}}}}

% ── Colours ───────────────────────────────────────────────────────────
\definecolor{gold}{HTML}{C9A84C}
\definecolor{darkbg}{HTML}{1A1A1A}

% ── Hyperlinks ────────────────────────────────────────────────────────
\hypersetup{
  colorlinks=true,
  linkcolor=black,
  urlcolor=gold,
  citecolor=black
}

% ── Header / Footer ──────────────────────────────────────────────────
\pagestyle{fancy}
\fancyhf{}
\setlength{\headheight}{14pt}
\addtolength{\topmargin}{-2pt}
\renewcommand{\headrulewidth}{0.4pt}
\fancyhead[L]{\small Agentic Finance System}
\fancyhead[R]{\small Software Engineering -- Semester 6}
\fancyfoot[C]{\thepage}

% ── Section formatting ───────────────────────────────────────────────
\titleformat{\section}{\Large\bfseries}{}{0em}{}[\vspace{-0.5em}\rule{\textwidth}{0.4pt}]
\titleformat{\subsection}{\large\bfseries}{}{0em}{}

% ══════════════════════════════════════════════════════════════════════
\begin{document}

% ── Title Page ────────────────────────────────────────────────────────
\begin{titlepage}
\centering
\vspace*{2cm}

{\Huge\bfseries Agentic Finance System}\\[0.6cm]
{\Large An AI-Powered Personal Finance and\\Market Insight Platform}\\[1.2cm]

\rule{0.6\textwidth}{0.4pt}\\[1.4cm]

{\large\textbf{Software Engineering Project Report}}\\[0.3cm]
{\large Semester 6}\\[2cm]

\begin{tabular}{l l}
\textbf{Jainithissh S}       & CB.SC.U4AIE23129  \\[4pt]
\textbf{Krishna Prakash S}   & CB.SC.U4AIE232138 \\[4pt]
\textbf{Nitheshkummar C}     & CB.SC.U4AIE23155  \\[4pt]
\textbf{Akhilesh Kumar S}    & CB.SC.U4AIE232170 \\[4pt]
\textbf{Adarsh Pradeep}      & CB.SC.U4AIE23109  \\
\end{tabular}

\vfill
{\large Amrita Vishwa Vidyapeetham}\\[0.3cm]
{\large \today}
\end{titlepage}

% ── Table of Contents ────────────────────────────────────────────────
\tableofcontents
\newpage

% ══════════════════════════════════════════════════════════════════════
%  1 — PROBLEM STATEMENT
% ══════════════════════════════════════════════════════════════════════
\section{Problem Statement}

\subsection{Context}
India's digital payments ecosystem generates millions of bank SMS and UPI notifications every day. Despite this wealth of personal financial data, most users have no automated way to convert these messages into actionable spending insights. Financial news arrives as a separate, disconnected stream with no link to an individual's actual money situation.

\subsection{User Pain Points}
\begin{enumerate}[leftmargin=1.5em]
  \item \textbf{Manual expense tracking:} Bank SMS and UPI notifications arrive daily, yet tracking spending remains an entirely manual process. Users must open each message and mentally categorise every transaction.
  \item \textbf{No personal risk quantification:} Users cannot measure their own financial risk --- EMI burden, savings rate, or the share of variable versus fixed expenses --- without spreadsheet-level effort.
  \item \textbf{Disconnected market news:} Financial news is presented generically. There is no mechanism to relate it to an individual's portfolio, income level, or spending behaviour.
  \item \textbf{Inaccessible advisory:} Professional financial advisors are out of reach for students and early-career earners. No affordable tool combines personal data, live market sentiment, and domain knowledge into a single, grounded recommendation.
\end{enumerate}

\subsection{Why Existing Solutions Fall Short}
\begin{itemize}[leftmargin=1.5em]
  \item \textbf{Expense trackers} (e.g.\ Walnut, Money Manager) rely on manual entry and provide no forward-looking intelligence.
  \item \textbf{Market apps} (e.g.\ Moneycontrol, ET Markets) deliver generic news with no personalisation.
  \item \textbf{Robo-advisors} (e.g.\ Groww, Kuvera) focus on mutual-fund selection and ignore day-to-day spending patterns.
  \item No existing product bridges \emph{personal transaction data}, \emph{real-time news sentiment}, and \emph{curated financial knowledge} into a single explainable system.
\end{itemize}

\subsection{Core Gap}
There is no single system that unifies personal transaction data, real-time news sentiment, and financial domain knowledge into agentic, explainable, personalised guidance. The \textbf{Agentic Finance System} addresses this gap.

\subsection{Objective}
Build an end-to-end AI-powered platform that:
\begin{enumerate}[leftmargin=1.5em]
  \item Ingests and structures personal finance data from bank SMS and UPI notifications automatically.
  \item Computes a personal risk and spending profile using ML classifiers.
  \item Understands financial news and market trends using NLP transformers.
  \item Generates grounded, explainable recommendations via a RAG-augmented multi-agent LLM pipeline.
  \item Teaches financial literacy through a Bayesian adaptive micro-learning engine.
\end{enumerate}

% ══════════════════════════════════════════════════════════════════════
%  2 — SYSTEM ARCHITECTURE
% ══════════════════════════════════════════════════════════════════════
\newpage
\section{System Architecture}

The system follows a \textbf{multi-agent architecture} with three autonomous agents coordinated through a FastAPI backend. Each agent operates on a \textbf{Two-Tier Pipeline}: Tier~1 provides fast, deterministic ML/rule-based inference; Tier~2 adds LLM-powered agentic reasoning via Groq (LLaMA~3.3 70B), with automatic fallback to Tier~1 if the LLM is unavailable.

\subsection{High-Level Layers}

\begin{table}[h!]
\centering
\begin{tabularx}{\textwidth}{l X}
\toprule
\textbf{Layer} & \textbf{Description} \\
\midrule
Client        & Android app (Kotlin, Jetpack Compose) with background \texttt{NotificationListenerService} for automatic SMS capture \\
API Gateway   & FastAPI backend (Python~3.11, Uvicorn) with Firebase JWT authentication \\
Agents        & Agent~A (Finance Analyzer), Agent~B (Market Intelligence), Agent~C (Decision Synthesizer) \\
Learning      & Bayesian Knowledge Tracing engine with AI-generated micro-learning cards \\
Persistence   & Supabase PostgreSQL (production) / SQLite (development) with env-var routing \\
External APIs & Groq Cloud LLM, Economic Times RSS feed, NSE live market data \\
\bottomrule
\end{tabularx}
\end{table}

\subsection{Agent Descriptions}

\textbf{Agent A --- Personal Finance Analyzer}
\begin{itemize}[leftmargin=1.5em]
  \item Parses bank SMS via a 3-tier regex extraction cascade (currency-pattern $\to$ keyword-proximity $\to$ digit fallback).
  \item Classifies transactions into 10 spending categories using TF-IDF + Logistic Regression.
  \item Evaluates financial risk with a Random Forest classifier (120 trees, depth~6) trained on 6~financial ratio features.
  \item Prorates partial-month data to prevent early-month low-risk bias.
  \item Generates a 4--5 sentence data-specific narrative via Groq LLM (Tier~2).
\end{itemize}

\textbf{Agent B --- News \& Market Intelligence}
\begin{itemize}[leftmargin=1.5em]
  \item Performs sentiment analysis on live Economic Times RSS headlines using FinBERT (ProsusAI/finbert).
  \item Extracts company and sector entities via spaCy NER with~25+ company-to-sector mappings.
  \item Predicts market trend (bullish / sideways / bearish) using a Random Forest on FinBERT aggregate features.
  \item Supplies real-time NSE index data pre-warmed at server startup.
\end{itemize}

\textbf{Agent C --- Decision Synthesizer \& RAG Chatbot}
\begin{itemize}[leftmargin=1.5em]
  \item Retrieves top-3 relevant chunks from a 32-chunk curated financial knowledge corpus via custom TF-IDF with tag-bonus ($\times 3$) weighting.
  \item Combines Agent~A financial data, Agent~B sentiment, and RAG context in a multi-turn Groq LLaMA~3.3 70B chatbot.
  \item Falls back to a keyword-matched deterministic response covering 6~core topics when the LLM is unavailable.
\end{itemize}

% ══════════════════════════════════════════════════════════════════════
%  3 — TECHNOLOGY STACK
% ══════════════════════════════════════════════════════════════════════
\newpage
\section{Technology Stack}

\begin{table}[h!]
\centering
\renewcommand{\arraystretch}{1.25}
\begin{tabularx}{\textwidth}{l l X}
\toprule
\textbf{Layer} & \textbf{Technology} & \textbf{Purpose} \\
\midrule
Backend       & FastAPI + Uvicorn (Python~3.11)            & Async REST API with JSON schema validation \\
Database      & Supabase PostgreSQL / SQLite               & Dual-DB with env-var routing and auto-failover \\
Auth          & Firebase Admin SDK + Google Sign-In        & JWT verification and per-user data isolation \\
ML --- Risk   & scikit-learn Random Forest                 & Financial risk classification (low / medium / high) \\
ML --- Category & TF-IDF + Logistic Regression             & 10-class SMS transaction categorisation \\
ML --- Trend  & Random Forest on FinBERT aggregates        & Market direction prediction \\
NLP --- Sentiment & HuggingFace FinBERT (ProsusAI)        & Finance-domain sentiment scoring \\
NLP --- NER   & spaCy \texttt{en\_core\_web\_sm}           & Company and sector entity extraction \\
LLM           & Groq Cloud (LLaMA~3.3 70B Versatile)      & Reasoning, chatbot, card generation \\
RAG Engine    & Custom TF-IDF + IDF smoothing              & Knowledge-grounded responses (32 chunks) \\
Mobile        & Kotlin + Jetpack Compose + Retrofit        & Android UI with background SMS capture \\
Deployment    & Docker + Render.com / HuggingFace Spaces   & Containerised cloud deployment \\
\bottomrule
\end{tabularx}
\end{table}

\subsection{Key Python Dependencies}
\texttt{fastapi}, \texttt{uvicorn}, \texttt{pydantic}, \texttt{scikit-learn}, \texttt{pandas}, \texttt{joblib}, \texttt{spacy}, \texttt{transformers} (FinBERT), \texttt{supabase}, \texttt{firebase-admin}, \texttt{groq}, \texttt{httpx}, \texttt{python-dotenv}.

\subsection{ML Models Summary}

\begin{table}[h!]
\centering
\renewcommand{\arraystretch}{1.2}
\small
\begin{tabularx}{\textwidth}{l l l X}
\toprule
\textbf{Model} & \textbf{Algorithm} & \textbf{Data} & \textbf{Output} \\
\midrule
Category   & TF-IDF + Log.\ Regression          & 1K+ synth.\ SMS        & 10 categories \\
Risk       & Random Forest (120, d=6)              & 900 synth.\ profiles   & low/med/high + conf.\ \\
Trend      & Random Forest (120, d=5)              & 600 synth.\ regimes    & bullish/sideways/bearish \\
FinBERT    & BERT (finance fine-tuned)             & ProsusAI corpus        & pos./neg./neutral \\
spaCy NER  & CNN + transition parser               & OntoNotes~5.0          & ORG/PERSON/LOC \\
Groq LLaMA & 70B autoregressive LLM               & Web-scale              & NL generation \\
\bottomrule
\end{tabularx}
\end{table}

% ══════════════════════════════════════════════════════════════════════
%  4 — FEATURES
% ══════════════════════════════════════════════════════════════════════
\newpage
\section{Features}

\subsection{Automatic SMS Transaction Parsing}
The system intercepts bank SMS and UPI notifications via Android's \texttt{NotificationListenerService}. Every notification is forwarded to the backend, which acts as the sole parsing authority. A 3-tier extraction pipeline handles the wide diversity of Indian bank message formats:
\begin{enumerate}[leftmargin=1.5em]
  \item \textbf{Tier~1 --- Currency Regex:} Matches patterns like \texttt{Rs}, \texttt{INR}, or the rupee symbol followed by an amount.
  \item \textbf{Tier~2 --- Keyword Proximity:} Searches for amounts within 15~characters of keywords ``debited'' or ``credited''.
  \item \textbf{Tier~3 --- Digit Fallback:} Extracts any standalone number with two or more digits as a last resort.
\end{enumerate}
Credit/debit classification uses 15 credit keywords versus 12 debit keywords. Category assignment combines a keyword map with a TF-IDF + Logistic Regression model at a confidence threshold $\geq 0.50$ across 10~classes.

\subsection{Personal Financial Risk Analysis (Agent A)}
Agent~A aggregates a user's transactions over a specified date range and computes a 6-feature vector: income, fixed expenses, variable expenses, savings, savings rate, and variable share. These feed into a Random Forest classifier that outputs a risk level (low, medium, high) with a confidence score. A month-proration step projects partial-month data to a full month, preventing the common bias of artificially low risk in the first few days of the month. The LLM tier generates a personalised 4--5 sentence narrative that references the user's actual rupee amounts.

\subsection{News \& Market Sentiment Intelligence (Agent B)}
Agent~B fetches live headlines from the Economic Times RSS feed and scores each article through FinBERT (headline weight 40\%, body weight 60\%). Thirteen boilerplate-removal regex patterns strip editorial noise before inference. spaCy NER extracts company entities, which are mapped to sectors (Banking, IT, Auto, Telecom, FMCG, Metals, Pharma) via a 25+ entry lookup. A Random Forest trend model aggregates per-article sentiments into an overall market direction. Real-time NSE index data is pre-warmed at server startup for zero cold-start latency.

\subsection{RAG-Grounded AI Chatbot (Agent C)}
Agent~C provides a multi-turn conversational interface. For every user query, it retrieves the top-3 most relevant chunks from a curated 32-chunk financial knowledge corpus (covering emergency funds, the 50/30/20 rule, debt management, SIP investing, tax saving, insurance, and more). The retrieval uses custom TF-IDF scoring with IDF smoothing and a tag-match bonus ($\times 3$). The retrieved context, combined with Agent~A and Agent~B outputs and full conversation history, is passed to Groq LLaMA~3.3 70B. The system prompt enforces rupee-denominated, data-specific answers. A deterministic keyword-matched fallback covers 6~core topics when the LLM is unavailable.

\subsection{Adaptive Micro-Learning System}
The learning engine uses \textbf{Bayesian Knowledge Tracing (BKT)} with a 3-state belief model per concept: $P(\text{Unknown}) + P(\text{Partial}) + P(\text{Mastered}) = 1.0$. Belief updates are conditioned on response correctness and latency:
\begin{itemize}[leftmargin=1.5em]
  \item Fast correct ($< 30$s): $P(\text{Mastered}) \mathrel{+}= 0.30$
  \item Slow correct ($> 30$s): $P(\text{Mastered}) \mathrel{+}= 0.15$, $P(\text{Partial}) \mathrel{+}= 0.15$
  \item Fast wrong: $P(\text{Unknown}) \mathrel{+}= 0.20$ (guessing interpretation)
  \item Slow wrong ($> 60$s): $P(\text{Unknown}) \mathrel{+}= 0.40$ (struggling interpretation)
\end{itemize}
A DAG-structured concept graph with 10~financial topics and 4~difficulty levels (Beginner through Expert) ensures prerequisite-based unlock. The curriculum compiler scores each candidate concept as $\text{score} = \text{readiness} \times \text{urgency} \times \text{relevance}$, surfacing the weakest unlocked concept most relevant to the user's financial profile. AI-generated learning cards (150--200 words + MCQ) are produced via Groq and cached per mastery level.

\subsection{Gamified Learning Roadmap}
Five lessons progress from Financial Basics through Investment Basics. A scoring system awards base points (50), perfect-score bonus (100), and speed bonus (20). Star ratings are assigned at 3-star ($\geq 90\%$), 2-star ($\geq 70\%$), and 1-star ($\geq 50\%$). Six achievements --- First Steps, Bookworm, Perfect Score, Speed Learner, Week Warrior, and Master --- provide additional engagement.

\subsection{Android Application}
The Android client comprises 30~Kotlin source files across 6~feature modules and 6~screens: Login (Firebase Google Sign-In), Home/Dashboard, Chat, Learn (vertical-swipe reels), Roadmap, and Market. The \texttt{NotificationListenerService} runs in the background, capturing all notification text and forwarding it to the backend. Kotlin Coroutines (\texttt{Dispatchers.IO}) ensure non-blocking HTTP calls via Retrofit~2 + OkHttp.

% ══════════════════════════════════════════════════════════════════════
%  5 — API REFERENCE
% ══════════════════════════════════════════════════════════════════════
\newpage
\section{API Reference}

The backend exposes 17~endpoints, all requiring Firebase Bearer JWT (development mode accepts unauthenticated requests as \texttt{default\_user}).

\begin{table}[h!]
\centering
\renewcommand{\arraystretch}{1.15}
\small
\begin{tabularx}{\textwidth}{l l l X}
\toprule
\textbf{Method} & \textbf{Endpoint} & \textbf{Agent} & \textbf{Description} \\
\midrule
POST   & /api/parse\_message               & ---   & Parse raw SMS into a structured transaction \\
GET    & /api/transactions                 & ---   & List user's recent transactions \\
DELETE & /api/transactions                 & ---   & Delete all transactions for user \\
PUT    & /api/transactions/\{id\}/category & ---   & Override a transaction category \\
POST   & /api/analyze\_finance             & A     & ML risk analysis + LLM insight \\
POST   & /api/expense\_breakdown           & A     & Category-wise spending breakdown \\
GET    & /api/risk\_logs\_summary           & A     & Historical risk evaluation summary \\
POST   & /api/analyze\_news                & B     & FinBERT topic-based sentiment analysis \\
GET    & /api/news\_feed                   & B     & Live RSS feed with FinBERT scores \\
GET    & /api/live\_market                 & B     & Real-time NSE indices and sectors \\
POST   & /api/synthesize                   & C     & RAG + multi-turn LLM chatbot \\
GET    & /api/learning/next-card           & Learn & Bayesian-selected micro-learning card \\
POST   & /api/learning/submit-answer       & Learn & Submit answer, trigger BKT update \\
GET    & /api/learning/roadmap             & Learn & Gamified roadmap data \\
POST   & /api/learning/complete-lesson     & Learn & Mark lesson complete, award XP \\
GET    & /api/learning/stats               & Learn & User stats (streak, XP, mastery) \\
POST   & /api/analyze\_article             & B     & Deep single-article analysis \\
\bottomrule
\end{tabularx}
\end{table}

% ══════════════════════════════════════════════════════════════════════
%  6 — DESIGN PATTERNS & SE PRINCIPLES
% ══════════════════════════════════════════════════════════════════════
\section{Design Patterns \& Software Engineering Principles}

\begin{table}[h!]
\centering
\renewcommand{\arraystretch}{1.2}
\begin{tabularx}{\textwidth}{l X}
\toprule
\textbf{Pattern} & \textbf{Application} \\
\midrule
Multi-Agent Architecture & Three autonomous agents with strict separation of concerns \\
Two-Tier Pipeline        & Traditional ML (Tier~1) + LLM Agentic (Tier~2) with fallback \\
Singleton                & GroqClient, ConceptGraph, FinBERT pipeline, Supabase client \\
Strategy                 & Dual-DB routing (SQLite vs Supabase) via environment variable \\
Observer                 & Android \texttt{NotificationListenerService} monitors notification stream \\
DAG Structure            & Concept prerequisite graph with DFS cycle detection \\
Bayesian Inference       & Probabilistic knowledge state tracking (not deterministic scores) \\
RAG Pattern              & TF-IDF retrieval + tag-bonus weighting grounds all LLM responses \\
\bottomrule
\end{tabularx}
\end{table}

Additional practices: CORS middleware, Firebase auth guards, background startup tasks (FinBERT pre-load in daemon thread), response caching, and graceful degradation via deterministic fallbacks at every LLM-dependent code path.

% ══════════════════════════════════════════════════════════════════════
%  7 — NOVELTY
% ══════════════════════════════════════════════════════════════════════
\newpage
\section{Novelty \& Key Contributions}

\subsection{Two-Tier Agentic Pipeline}
Unlike conventional systems that rely solely on either ML models or LLMs, every agent in the Agentic Finance System operates on a \textbf{Two-Tier Pipeline}. Tier~1 provides fast, deterministic ML/rule-based inference that is always available. Tier~2 layers LLM-powered agentic reasoning on top. If the LLM fails or experiences latency spikes, Tier~1 automatically takes over. This architecture guarantees \textbf{100\% core feature availability} regardless of external API status --- a critical requirement for a financial application.

\subsection{Unified Personal + Market Context}
No existing consumer product unifies a user's personal transaction data, real-time news sentiment, and curated financial knowledge into a single agentic response. Agent~C of this system does exactly that: it injects Agent~A's financial profile, Agent~B's sentiment signals, and RAG-retrieved domain knowledge into every LLM call, producing advice that references the user's actual rupee amounts and current market conditions.

\subsection{Bayesian Adaptive Learning for Finance}
The micro-learning engine applies Bayesian Knowledge Tracing --- a technique proven in educational technology --- to personal finance education. The 3-state belief model (Unknown, Partial, Mastered) with latency-aware updates provides a more nuanced understanding of learner knowledge than binary correct/incorrect scoring. The curriculum compiler dynamically prioritises concepts based on readiness, urgency, and relevance to the user's financial profile (e.g.\ a high-risk user is steered toward budgeting content first).

\subsection{Backend-Authoritative SMS Parsing}
The Android client sends \emph{all} notifications to the backend without any client-side filtering. This design choice places the parsing intelligence entirely server-side, enabling:
\begin{itemize}[leftmargin=1.5em]
  \item Rapid iteration on parsing rules without app updates.
  \item A single source of truth for transaction classification.
  \item Consistent behaviour across all Android devices and OS versions.
\end{itemize}

\subsection{Domain-Specific NLP Stack}
Rather than using a general-purpose sentiment model, the system employs \textbf{FinBERT} --- a BERT model fine-tuned specifically on financial text. Combined with 13~boilerplate-removal patterns, 27~positive and 31~negative domain keyword fallbacks, and strong-phrase detection, the pipeline achieves finance-aware sentiment accuracy that generic models cannot match.

\subsection{Graceful Degradation by Design}
Every feature that depends on an external API (Groq LLM, Economic Times RSS, NSE data) has a deterministic fallback path. The system never returns an empty or error state to the user; instead, it degrades gracefully to rule-based or cached responses. This is a deliberate architectural decision, not an afterthought.

% ══════════════════════════════════════════════════════════════════════
%  8 — FUTURE SCOPE
% ══════════════════════════════════════════════════════════════════════
\newpage
\section{Future Scope}

\subsection{AI \& Model Enhancements}
\begin{itemize}[leftmargin=1.5em]
  \item Replace custom TF-IDF RAG with FAISS or Pinecone vector store for semantic similarity retrieval.
  \item Add a self-evaluation loop where Agent~C critiques and refines its own output before delivery.
  \item Incorporate reinforcement learning from explicit and implicit user feedback.
  \item Introduce LSTM or Prophet time-series models for predictive spending forecasting.
\end{itemize}

\subsection{Product Features}
\begin{itemize}[leftmargin=1.5em]
  \item Voice-based financial assistant via speech-to-text integration.
  \item Portfolio tracking with Zerodha or Groww brokerage API connections.
  \item Multi-language SMS parsing for Hindi and regional Indian languages.
  \item Play Store release and cross-platform iOS support via React Native.
\end{itemize}

\subsection{Engineering \& Operations}
\begin{itemize}[leftmargin=1.5em]
  \item MLflow experiment tracking and model registry for version management.
  \item Formal evaluation framework: accuracy, AUC-ROC, and human judgement scores.
  \item CI/CD pipeline with GitHub Actions for automated testing and deployment.
  \item On-device SMS parsing option for enhanced user privacy.
\end{itemize}

% ══════════════════════════════════════════════════════════════════════
%  CONCLUSION
% ══════════════════════════════════════════════════════════════════════
\section{Conclusion}

The Agentic Finance System demonstrates that a multi-agent, two-tier architecture can deliver personalised, grounded financial guidance by unifying personal transaction data, real-time market sentiment, and curated domain knowledge. With 3~AI agents, 17~API endpoints, 6~ML/AI models, a 32-chunk RAG corpus, Bayesian adaptive learning across 10~financial concepts, and a full Android client, the system covers the end-to-end pipeline from raw SMS ingestion to explainable financial advice. The deliberate emphasis on graceful degradation ensures that every feature remains operational even when external services are unavailable --- a principle essential for any production financial application.

\end{document}
